\section*{(2a aula) Espaços Vetoriais (seções 1.1-1.3)}

A ferramenta matemática essencial da Mecânica Clássica é a Teoria de Equações Diferenciais Ordinárias. 
Na Mecânica Quântica (MQ) temos a \textbf{Álgebra Linear} e a \textbf{Teoria das Equações Diferenciais Parciais} 
(EDP) como fundamentais. Nosso curso vai começar com uma revisão de \textbf{Álgebra Linear}. As EDP's aparecerão 
mais tarde.

Por que Álgebra Linear? Como veremos, nos postulados da MQ, o estado de uma partícula é representado, não por 
uma função vetorial ${\bm r}(t)$ como na Mecânica Clássica, mas sim como um elemento de um espaço vetorial, 
representado pelo símbolo $\ket{\psi(t)}$. A conexão desses objetos com a função de onda probabilística 
$\Psi\prt{\mathbf{r},t}$ será discutida quando estabelecermos os postulados. Esqueça a Física por enquanto 
e se concentre em entender as propriedades dos \textbf{ESPAÇOS VETORIAIS}.
\begin{shadequote}
    Um \textbf{ESPAÇO VETORIAL} $V$ é uma coleção de objetos chamados \textbf{VETORES} 
    $\ket{v_1}$, $\ket{v_2}$, $\ket{v_3}$, $\dots$ com duas operações definidas
    \begin{align}
        \text{Soma :} \quad &  \ket{u} + \ket{v} \in V \\
        \text{Multiplicação por escalar :} \quad & a\ket{v} \in V 
    \end{align}
\end{shadequote}


Esse escalar pode ser um número real ou complexo, no primeiro temos um \textbf{E.V. REAL} no outro \textbf{E.V. COMPLEXO}. 
\textbf{Nota}: O vetor não é nem real nem complexo! As propriedades da soma de vetores e multiplicação por escalares 
são as que você esperaria:

\begin{enumerate}
    \item $a(\ket{u} + \ket{w}) = a\ket{u} + a\ket{w}$ \quad (distributiva)
    \item $(a + b)\ket{v} = a\ket{v} + b\ket{v}$ \quad (distributiva)
    \item $a(b\ket{v}) = (ab)\ket{v}$ \quad (associativa)
    \item $\ket{v} + \ket{w} = \ket{w} + \ket{v}$ \quad (comutativa)
    \item $\ket{v} + (\ket{w} + \ket{z}) = (\ket{v} + \ket{w}) + \ket{z}$ \quad (associativa)
    \item $\ket{v} + \ket{0} = \ket{v}$ \quad (existe o vetor nulo)
    \item $\ket{v} + (-\ket{v}) = \ket{0}$ \quad (existe o inverso do vetor)
\end{enumerate}

É um jogo que os matemáticos adoram demonstrar outras propriedades dos espaços vetoriais usando apenas as 
propriedades acima.



{\bf Exemplo: } Prove que só existe um único vetor nulo em $V$.
Suponha que $\ket{0}$ e $\ket{0'}$ sejam vetores nulos.
\begin{align*}
    \ket{0} + \ket{0'} &= \ket{0} \quad \text{(por 6, pois $\ket{0'}$ é nulo)} \\
    \ket{0} + \ket{0'} &= \ket{0'} \quad \text{(por 6, pois $\ket{0}$ é nulo)}
\end{align*}
Logo, $\ket{0} = \ket{0'}$ (são os mesmo vetores).

{\bf Exemplo :} Setas em um plano formam um espaço vetorial real. 
Aqui você tem que informar como se somam setas e como se multiplicam 
setar por escalares.


\begin{figure}[htb]
    \centering
    \begin{subfigure}[b]{0.5\textwidth}
        %\documentclass[tikz,multi=true]{standalone}
%\usetikzlibrary{arrows.meta}
%\begin{document}
\begin{tikzpicture}[>=Stealth, scale=1.]
% Origem
\coordinate (O) at (0,0);

% Vetores V e W (modifique as coordenadas conforme desejar)
\coordinate (V) at (5,1);    % Vetor V
\coordinate (W) at (1,2);    % Vetor W
\coordinate (S) at ($(V)+(W)$);  % Soma V + W

% Desenha os vetores V e W
\draw[->, thick, black] (O) -- (V) node[midway, below right]{$\vec{V}$};
\draw[->, thick, black] (V) -- (7,5) node[midway, below right]{$\vec{W}$};
\draw[->, thick, black] (O) -- (7,5) node[midway, above left]{$\vec{V} + \vec{W}$};

\end{tikzpicture}
%\end{document}


    \end{subfigure}
    \begin{subfigure}[b]{0.4\textwidth}
        %\documentclass[tikz,multi=true]{standalone}
% \usetikzlibrary{arrows.meta}
% \begin{document}
\begin{tikzpicture}[>=Stealth, scale=1.]

\draw[->, thick] (0,0.1) -- (7,0.1) node[midway ,above]{$a\vec{V}$};
\draw[->, thick] (0,0) -- (5,0) node[midway ,below]{$\vec{V}$};
\end{tikzpicture}

%\end{document}
        \vspace{1cm}
        % \documentclass[tikz,multi=true]{standalone}
% \usetikzlibrary{arrows.meta}

%\begin{document}

    \begin{tikzpicture}[>=Stealth, scale=1.0]
    
        \draw[->, thick] (0,0.0) -- (3,0.0) node[midway ,above]{$\vec{V}$};
        \draw[->, thick] (0,0.0) -- (-3,0.0) node[midway ,above]{$\vec{W}$};
        \filldraw[black] (0,0) circle (1.5pt)
    
    \end{tikzpicture}

%\end{document}
        \vspace{1cm}
    \end{subfigure}
    \caption{é a maneira de somar $V$ e $W$. $V+W$ é uma seta no plano.
    é uma maneira de multiplicar por um numero real (se $a<0$ inverta a seta).$W$ é o inverso de $V$}
\end{figure}
${\vec 0}$ é uma seta de tamanho zero.

{\bf Exemplo :} O conjunto de todas as matrizes reais $2\times2$ formam 
um espaço vetorial real se as operações de soma e multiplicação por escalar forem, respectivamente

\[
    \begin{bmatrix}
    a & b \\
    c & d
    \end{bmatrix}
    +
    \begin{bmatrix}
    a' & b' \\
    c' & d'
    \end{bmatrix}
    =
    \begin{bmatrix}
    a + a' & b + b' \\
    c + c' & d + d'
    \end{bmatrix}
    \qquad \text{e} \qquad  
    \alpha \begin{bmatrix}
    a & b \\
    c & d
    \end{bmatrix}
    =
    \begin{bmatrix}
    \alpha a & \alpha b \\
    \alpha c & \alpha d
    \end{bmatrix}
\]

Com essas definições:
\begin{itemize}
    \item Vetor nulo: $\begin{bmatrix} 0 & 0 \\ 0 & 0 \end{bmatrix}$
    \item Inverso de $\begin{bmatrix} a & b \\ c & d \end{bmatrix}$ é $\begin{bmatrix} -a & -b \\ -c & -d \end{bmatrix}$
\end{itemize}

No entanto, o conjunto das matrizes da forma $\begin{bmatrix} a & b \\ c & 1 \end{bmatrix}$  (com $a, b, c$ reais) 
{\bf não} forma um espaço vetorial com essas regras, pois:
\[
\begin{bmatrix} a & b \\ 1 & c \end{bmatrix}
+
\begin{bmatrix} a' & b' \\ 1 & c' \end{bmatrix}
=
\begin{bmatrix} a + a' & b + b' \\ c+c' & 2 \end{bmatrix}
\]
não pertence ao conjunto.
\begin{observação}
Não se preocupe com as 7 regras acima, elas impõem propriedades 
inteiramente naturais para a soma de vetores e multiplicação pode escalar.
Algumas propriedades derivadas dessas 7, que à  principio deveriam ser
demonstradas, podem ser lembradas por referencia ao espaco vetorial com
que você está mais familiarizado, o conjunto de setas no plano. Por exemplo
(tente demonstrar) : 

\begin{equation*}
    0\ket{0} = \ket{0} & \longrightarrow \underbrace{(0+1)\ket{V}}_{\text{(por 2)}} =
        0\ket{V} + 1\ket{V} & \longrightarrow 1\ket{V} = 0\ket{V} + 1\ket{V} = 0\ket{0} \longrightarrow 0\ket{0} = \ket{0} \\
\end{equation*}
%     0\ket{0} = \ket{0} \longrightarrow (0+1)\ket{V} = \text{(por 2)} \\ 
%     0\ket{V} + 1\ket{V} \longrightarrow 1\ket{V} = 0\ket{V} + 1\ket{V} = 0\ket{0} \longrightarrow 0\ket{0} = \ket{0} \\
% \]
\begin{equation}
    \underbrace{\ket{-V}}_{\text{(o inverso de $\ket{V}$ é unico)}} = \ket{0} =(por *) (1-1)\ket{V} = (por 2) ->  1\ket{V} + (-1)\ket{V}
    como 1\ket{V} = \ket{V}, segue de 7 que -\ket{V} = \ket{-V}
\end{equation}

\begin{shadequote}
Um conjunto de $N$ vetores não-nulos é \textbf{LINEARMENTE INDEPENDENTE} se a equação:
    \begin{equation}
        \sum_{i=1}^N a_i \ket{v_i} = \ket{0}    
    \end{equation}
só tem a solução trivial (com todos $a_i = 0$). Caso exista uma solução não-trivial, o conjunto é {\bf LINEARMENTE DEPENDENTE}.
\end{shadequote}



{\bf Exemplo :} Dois vetores no plano, não colineares, são LI. Três vetores no plano são LD, pois um deles pode 
ser escrito como combinação dos outros dois.(FAZER IMAGENS). Esse exemplo, por se tratar de um espaço vetorial 
de naturezageométrica, não ilustra bem o processo usado para se determinar se um conjunto de vetores é LI ou LD.

{\bf Exemplo :}
$\ket{1} = \begin{bmatrix}
    0 & 1 \\
    0 & 0
    \end{bmatrix}$, 
    $\ket{2} = \begin{bmatrix}
    1 & 1 \\
    0 & 1
    \end{bmatrix}$
$\ket{3} = \begin{bmatrix}
    -2 & -1 \\
    0 & -2
    \end{bmatrix}$
são LI ou LD ? 
formamos uma combina;cao linear generica e igualamos a $\ket{0}$
$$
a\ket{1}+b\ket{2}+c\ket{3} = (pelas regras de soma e multiplicaáo por escalar)\begin{bmatrix}
    b-2c & a+b+c \\
    0 & b-2c
    \end{bmatrix}
    =
    \begin{bmatrix}
    0 & 0 \\
    0 & 0
    \end{bmatrix} =\ket{0}
$$
a soluçao dessa equaçao vetorial é obtida resolvendo 
\begin{equation}
    \begin{array}{l}
        b - 2c = 0 \\
        a + b -c = 0 \\
        b - 2c = 0
    \end{array} 
\end{equation}
claramente $a=-1$, $b=2$ e $c=1$ é uma solução não trivial  logo vetores LD. 


\begin{shadequote}
    A \textbf{DIMENSÃO} de um espaço vetorial (EV) é o número máximo de vetores LI que ele pode 
    acomodar. $\mathbb{V}^n(R)$ é um espaço vetorial real de dimensão $n$, e $\mathbb{V}^n(C)$ 
    é um  espaço vetorial complexo de dimensão $n$.
\end{shadequote}

{\bf Exemplo:} O EV de setas no plano tem dimensao 2 pos duas setas não colineares são LI e 
setas quaisquer são LD.

{\bf Exemplo:} O EV matrizes reais $2\times2$ tem dimensão 4 pois
\begin{equation}
    \begin{bmatrix}
    1 & 0 \\
    0 & 0
    \end{bmatrix},
    \begin{bmatrix}
    0 & 1 \\
    0 & 0
    \end{bmatrix},
    \begin{bmatrix}
    0 & 0 \\
    1 & 0
    \end{bmatrix},
    \begin{bmatrix}
    0 & 0 \\
    0 & 1
    \end{bmatrix}
\end{equation}
são LI. (mostre isso) e qualquer outro vetor $\begin{bmatrix}
    0 & 0 \\
    0 & 1
    \end{bmatrix}$ pode ser escrito como combinação linear desses 4 vetores. 

Você percebe que, para provar que a dimensão de um EV é $n$ você deve : 
\begin{itemize}
    \item exibir $n$ vetores LI
    \item mostrar que um vetor genérico pode ser escrito como uma combinação linear 
    desses $n$ vetores. 
\end{itemize}

{\bf TEOREMA:} Qualquer vetor em um espaço de dimensão $n$ pode ser escrito como combinação
linear de $n$ vetores LI. 
    
\begin{equation}
    \ket{V} = \sum^n_{i=1}{v_i}\ket{i}
\end{equation}
PROVA : Suponha que exista um $\ket{V}$ para o qual não seja possível encontrar esses 
coeficientes $v_i$, isso implicaria que $\ket{V}$, mais os $n$ vetores $\ket{i}$, são
LI o que não pode acontecer em um espaço de dimensão $n$.

{\bf DEFINIÇÃO :} Qualquer conjunto de $n$ vetores LI em um $\mathbb{V}^n$ forma uma BASE.
{\bf TEOREMA :} As componentes de um vetor em uma base são únicos
{\bf PROVA :}
\begin{equation}
    \ket{V} = \sum^n_{i=1}{v_i}\ket{i}
\end{equation}
\begin{equation}
    \ket{V} = \sum^n_{i=1}{v^\prime}\ket{i}
\end{equation}
subtraindo
\begin{equation}
    \ket{0} = \sum^n_{i=1}{v_i-v^\prime_i}\ket{i}
\end{equation}
como a base é necessariamente LI -> $v_i=v^\prime_i$ (pela definição de vetores LI)


As equações da Mec. Clássica e do Eletromagnetismo são equações que relacionam vetores do tipo
"seta". 
\begin{equation*}
\begin{split}
    &\vec{F} = m\vec{a}, \\ 
    \nabla\times&\vec{E}+\frac{1}{c} \frac{\partial \vec{B}}{\partial t} = 0, \ldots
\end{split}
\end{equation*}
Essas setas $\vec{F},\vec{B},\vec{E},$ etc existem independentemente de uma base, no entanto pra fazer
cálculos escolhemos um sistema de eixos cartesianos com 3 versores $\hat{i},\hat{j},\hat{k}$. Em termos
desta base os vetores das equações são expressos a cada equaçào vetorial resultam em 3 equações
diferenciais para as componentes dos vetores. Por exemplo : 

\begin{equation}
\nabla\times&\vec{E}+\frac{1}{c} \frac{\partial \vec{B}}{\partial t} = 0 \Longrightarrow \left\{\begin{array}{l}
    \dfrac{\partial E_z}{\partial y} - \dfrac{\partial E_y}{\partial z} + \dfrac{1}{c} \dfrac{\partial B_x}{\partial t} = 0 \\
    \dfrac{\partial E_x}{\partial z} - \dfrac{\partial E_z}{\partial x} + \dfrac{1}{c} \dfrac{\partial B_y}{\partial t} = 0 \\
    \dfrac{\partial E_y}{\partial x} - \dfrac{\partial E_x}{\partial y} + \dfrac{1}{c} \dfrac{\partial B_z}{\partial t} = 0
\end{array}\right.
\end{equation}

As equações da Mecânica Quantica são relações entre vetores de um espaço vetorial de estados físicos. Os  cálculos 
são feitos (e as EPS da MQ são obtidas) quando escolhemos uma base para o espaço vetorial e escrevemos os vetores
nessa base.

\hline
Vimos na Mecanica Clássica a relcào existente entre as componentes de um mesmo vetor $v$ em duas bases distintas.
FAZER DESENHO
${\bm v} = v_x\hat{i} + v_y\hat{j} + v_z\hat{k} = v_1\hat{e}_1 + v_2\hat{e}_2 + v_3\hat{e}_3$
\begin{equation*}
    \begin{bmatrix} v_x \\ v_y \\ v_z \end{bmatrix} =\begin{bmatrix} \bigLambda \end{bmatrix}
    \begin{bmatrix} v_1 \\ v_2 \\ v_3 \end{bmatrix}
\end{equation*}
\hline

As operações de norma de vetores e multiplicação por escalar para vetores escritos em uma mesma base, fica 

\begin{equation}
    \ket{V} = \sum^n_{i=1}{v_i}\ket{i}     \ket{W} = \sum^n_{i=1}{w_i}\ket{i}
\end{equation}
\begin{equation}
    \ket{V} + \ket{W} = \sum^n_{i=1}{v_i+w_i}\ket{i}
\end{equation}
\begin{equation}
    \alpha\ket{V} = \sum^n_{i=1}{\alpha v_i}\ket{i}
\end{equation}


\subsection*{Produto Interno}

O produto interno associa um número (real ou complexo) a um par de vetores. No espaço vetorial das "setas" existe 
uma definição geométrica para o produto interno:
\[
{\bm v} \cdot {\bm w} = |{\bm v}| |{\bm w}| \cos \theta
\]

onde $|{\bm v}|$ é o tamanho da seta ${\bm v}$ e o $\theta$ é o ângulo entre ${\bm v}$ e ${\bm w}$. No caso de 
EV gerais exigimos 3 propriedades de qualquer produto interno a ser proposto.
\begin{enumerate}
    \item $\bra{U}\ket{W} = \bra{W}\ket{U}^*$
    \item $\bra{V}\ket{V} \geq 0$, e $0$ apenas se $\ket{V} = \ket{0}$
    \item $\bra{V}[a\ket{W} + b\ket{Z}] = a\bra{V}\ket{W} + b \bra{V}\ket{Z}$
\end{enumerate}
aqui denotamos o produto interno pelo simbolo $\bra{U}\ket{W}$ e consideramos o caso geral onde 
esse número é complexo (caso estejamos analisando um produto interno real a conjugação em 1. não
tem efeito). O produto interno de setas obedece essas 3 propriedades. 

Em um EV generico se pode definir um produto interno por meio dos 
valiores atribuidos aos produtos internos de uma base. Veja 

\begin{equation}
    Se \ket{V} = \sum^n_{i=1}{v_i}\ket{i}  e   \ket{W} = \sum^n_{i=1}{w_i}\ket{i}
\end{equation}
\begin{equation}
    \braket{V}{W} = \sum^n_{i=1}{w_i} \bra{v}\ket{i} por (3)
\end{equation}
\begin{equation}
    \braket{V}{i} = \braket{i}{V}^* por (1)
\end{equation}
\begin{equation}
    \braket{i}{V}^* = \cch{\sum_{j=1}^n v_j \braket{i}{j}}^*  por (3)
\end{equation}
\begin{equation}
     = {\sum_{j=1}^n v^*_j \braket{i}{j}}  (conjugaçao complexa de unm produto)
\end{equation}
\begin{equation}
     = {\sum_{j=1}^n v^*_j \braket{j}{i}}  por (1)
\end{equation}
e finalmente
\begin{equation}
     \braket{V}{W} = {\sum_{i=1}^n\sum_{j=1}^n v^*_j w_i\bra{j}\ket{i}} 
\end{equation}
se conhecermos os valores de $\braket{i}{j}$, os protudos internos da base, podemos calcular 
qualquer produto interno.

\begin{shadequote}
    Uma base é ORTONORMAL quando o produto interno de seus elementos satisfaz $\braket{i}{j}=\delta_{ij}$
\end{shadequote}
 
 Em uma base ortonormal, temos:
\[
\braket{V}{W} = \sum_{i=1}^n v_i^* w_i
\]
veja que essa regra respeita as propriedades 1-3, em particular 
\[
\braket{V}{V} = \sum_{i=1}^n |v_i|^2 \geq 0
\]
Definimos como a norma de um vetor
\[
    |v_i| = \sqrt{\braket{V}{V}}
\]
onde $\braket{V}{V}$ é um número real por (1)

{\bf EXEMPLO :} No caso das setas, se usamos como base um conjunto de setar de tamanho 1 
e mutuamente ortogonais
\[
    \hat{i}\cdot\hat{i} = \hat{j}\cdot\hat{j} = \hat{k}\cdot\hat{k} = 1
\]
\[
    \hat{i}\cdot\hat{j} = \hat{i}\cdot\hat{k} = \hat{j}\cdot\hat{k} = 0
\]
o produto interno de $v=v_1\hat{i}+v_2\hat{j}+v_3\hat{k}$ e $w=w_1\hat{i}+w_2\hat{j}+w_3\hat{k}$ 
fica 
\begin{equation}
    {\bm v}\cdot{\bm w} = \sum_{i=1}^3 v_i\cdot w_i
\end{equation}

{\bf EXEMPLO :} No caso das matrizes $2\times2$ podemos usar como base
\[
    \ket{1}=\begin{bmatrix}
    1 & 0 \\
    0 & 0
    \end{bmatrix},
    \ket{2}=\begin{bmatrix}
    0 & 1 \\
    0 & 0
    \end{bmatrix},
    \ket{3}=\begin{bmatrix}
    0 & 0 \\
    1 & 0
    \end{bmatrix},
    \ket{4}=\begin{bmatrix}
    0 & 0 \\
    0 & 1
    \end{bmatrix}
\]
e definir um produto interno por meio de $\braket{i}{j}=\delta_{ij}$. Assim 
\[
    \ket{V}=v_1\ket{1}+v_2\ket{2}+v_3\ket{3}+v_4\ket{4} = 
    \begin{bmatrix}
        v_1 & v_2 \\
        v_3 & v_4
    \end{bmatrix} 
\]
e 
\[
    \braket{V}{W} = \sum_{i=1}^4v_iw_i
\]
A norma de $\ket{V} = 
    \begin{bmatrix}
        1 & 2 \\
        0 & 1
    \end{bmatrix}$ é 

\[
|V| = \sqrt{\braket{V}{V}} = \sqrt{1^2+2^2+0^2+1^2} = \sqrt{6}
\]
OBS : Ao definir um produto interno por meio dos produtos internos da base deve-se ter
cuidado para que esses respeitem as propriedades (1) e (2). No caso de matrizes não
poderíamos invertar um produto interno com os valores
\[
    \braket{1}{1}=-1 ou \braket{1}{2} = 1 e \braket{2}{1} = 0
\]
pois estaríamos violando 2 no primeiro caso e 1 no segundo caso.

OBS: No caso de uma base NÃO ORTOGONAL o produto interno deve ser ser calculado usando 
poderíamos inventar um produto interno com os valores
\[
    \braket{V}{W} = \sum_{i=1}^{n}\sum_{j=1}^{n}v^*_jw_i \braket{j}{i}
\]
e você deve fornecer os valores de $\braket{j}{i}$

\hline
    Essa expressao pode ser entendida como o produto matricial abaixo
\[
\begin{bmatrix}
v_1^* & v_2^* & \cdots & v_n^*
\end{bmatrix}
\begin{bmatrix}
\langle 1|1 \rangle & \langle 1|2 \rangle & \cdots & \langle 1|n \rangle \\
\langle 2|1 \rangle & \langle 2|2 \rangle & \cdots & \langle 2|n \rangle \\
\vdots & \vdots & \ddots & \vdots \\
\langle n|1 \rangle & \langle n|2 \rangle & \cdots & \langle n|n \rangle
\end{bmatrix}
\begin{bmatrix}
w_1 \\
w_2 \\
\vdots \\
w_n
\end{bmatrix}
= \langle V | W \rangle
\]
\hline
Os produtos internos da base são num certo sentido, o que define o produto interno.

\subsection*{Espaço Vetorial Dual}
Essa diferença no forma camo o vetor da esquerda é tratado na expressão do produto 
interno (seus correspondentes na base são conjugadas) nos leva a definir um EV dual as 
nosso EV original. Os vetores desse EVE deual usam o símbolo $\bra{V}$ (chamados bra)
para distingui-los dos vetores $\ket{V}$ (chamados ket) do EV original. Existe uma 
correspondencia entre os $\ket{V}$ do EV original e seus irmãos $\bra{V}$ do EV dual.
Se $\ket{U}+\ket{Z}=\ket{W}$, então no espaço dual devemos ter $\bra{U}+\bra{Z}=\bra{W}$.
Se $\ket{aU}=a\ket{U}$, então no espaço dual devemos ter $\bra{aU}=\bra{U}a^*$

Porque essa conjugaçao do escalar ? Para que a regra (1) e a regra(3) do produto interno
sejam respeitadas. O ponto é que esse espaço dual é apenas uma maneira de lembrar das 
regras do produto interno com facilidade. O símbolo $\braket{W}{U}$ ao invés de ser 
lido como  "o produto interno de $\ket{U}$ com $\ket{W}$, com $\ket{W}$ estando 
à esquerda" fica "o produto interno do bra $\bra{W}$ com o ket $\ket{U}$". Se $\ket{W}=\sum_{i=1}^nw_i\ket{i}$, então o bra associado fica $\bra{W}=\sum_{i=1}^nw^*_i\bra{i}$
e o produto interno pode ser calculado facilmente
\[
   \underbrace{\left(\sum_{i=1}^n w^*_i \bra{i}\right)}_{\bra{W}} \underbrace{\left(\sum_{j=1}^n u_j \ket{j}\right)}_{\ket{U}} = \sum_{i,j=1}^n w_i^* u_j \braket{i}{j}
\]
exatamente como deveria.

A operação de achar o bra correspondente a um determinado ket se chama OPERAÇÃO ADJUNTA. Pense no 
EV dual como uma espécia de imagem no espelho do EV original, onde cada ket tem seu 
correspondente (com o pequeno detalhe da conjugação dos escalares). Essa imagem ajuda a 
operar com o produto interno

\begin{figure}[!htbp]
\centering
    \includegraphics[width=0.5\linewidth]{class2/figs/fig4/dual_space.pdf}
    \caption{a Esquerda o espaço vetorial dual. A direita o espaço vetorial original}
    \label{fig:enter-label}
\end{figure}

O produto interno permite exprimir as components de um vetor em uma base na forma de um 
produto interno
\[
\ket{U} = \sum_{i=1}^n u_i \bra{i}
\]
Se a BASE FOR ORTONORMAL, podemos tornar o produto interno com $\bra{j}$ (dual de um dos 
elementos da base).
\[
\braket{j}{U} = \sum_{i=1}^n u_i\underbrace{\braket{j}{i}}_{\delta_{ij}}=U_j
\]
então a componente $u_j$ é a projeção de $\ket{U}$ na "direção" $\ket{j}$ (veja 
a lista de exercícios para o caso da base não ortonormal). Isso é análogo ao que 
acontece com setas (FAZER DESENHO)
\[
{\bm V} = v_i\hat{i}+v_2\hat{j} \Longrightarrow {v_1} = V\cdot\hat{i}\quad e \quad v_2 = V\cdot\hat{j}
\]

OBS: A única comunicaçãoentre o EV original e o EV dual é no produto interno, quando 
um bra e um ket produzem um número. NÃO FAZ SENTIDO SOMAR UM BRA E UM KET ! $\bra{U}+\ket{W}$ NÃO
EXISTE.

\end{document}